\begin{abstract}
  Kepler's first two laws state that a planet moves on an ellipse with the Sun at a focus and sweeps out equal areas in equal times (constant areal speed). In the \emph{Principia}, Newton showed how these laws connect to universal gravitation. Since then, the equivalence between orbital laws and force laws has remained a central topic in celestial mechanics. We present fully geometric proofs---built from explicit Euclidean straightedge-and-compass constructions---of this equivalence in both directions. One contribution is to use the auxiliary circle as the primary hodograph proxy in configuration space (rather than the directrix-circle normalization of radius \(2a\)), which streamlines several forward and inverse geometric constructions. Our emphasis is a \emph{Principia}-style argument that avoids differential equations while remaining close to Euclidean methods. We hope these proofs provide geometric insights complementary to Newton's analyses.
\end{abstract}
